\documentclass[10pt]{article}
\usepackage[utf8]{inputenc}
\usepackage[T1]{fontenc}
\usepackage{amsmath}
\usepackage{amsfonts}
\usepackage{amssymb}
\usepackage[version=4]{mhchem}
\usepackage{stmaryrd}
\usepackage{bbold}
\usepackage{hyperref}
\hypersetup{colorlinks=true, linkcolor=blue, filecolor=magenta, urlcolor=cyan,}
\urlstyle{same}

\title{Constructing new APN functions from known ones }

\author{Lilya Budaghyan ${ }^{\mathrm{a}, *}$, Claude Carlet ${ }^{\mathrm{b}}$, Gregor Leander ${ }^{\mathrm{c}}$\\
${ }^{\mathrm{a}}$ Department of Informatics, University of Bergen, PB. 7803, Bergen, Norway\\
${ }^{\mathrm{b}}$ Department of Mathematics, University of Paris 8 (MAATICAH), France\\
${ }^{\mathrm{c}}$ Department of Mathematics, Technical University of Copenhagen, Denmark}
\date{}


%New command to display footnote whose markers will always be hidden
\let\svthefootnote\thefootnote
\newcommand\blfootnotetext[1]{%
  \let\thefootnote\relax\footnote{#1}%
  \addtocounter{footnote}{-1}%
  \let\thefootnote\svthefootnote%
}

%Overriding the \footnotetext command to hide the marker if its value is `0`
\let\svfootnotetext\footnotetext
\renewcommand\footnotetext[2][?]{%
  \if\relax#1\relax%
    \ifnum\value{footnote}=0\blfootnotetext{#2}\else\svfootnotetext{#2}\fi%
  \else%
    \if?#1\ifnum\value{footnote}=0\blfootnotetext{#2}\else\svfootnotetext{#2}\fi%
    \else\svfootnotetext[#1]{#2}\fi%
  \fi
}

\begin{document}
\maketitle


\section*{ARTICLE INFO}
\section*{Article history:}
Received 22 May 2007\\
Revised 24 October 2008\\
Available online 13 November 2008 Communicated by Gary L. Mullen

\section*{Keywords:}
Almost bent\\
Almost perfect nonlinear\\
CCZ-equivalence\\
Differential uniformity\\
Nonlinearity\\
S-box\\
Vectorial Boolean function

\begin{abstract}
We present a method for constructing new quadratic APN functions from known ones. Applying this method to the Gold power functions we construct an APN function $x^{3}+\operatorname{tr}\left(x^{9}\right)$ over $\mathbb{F}_{2^{n}}$. It is proven that for $n \geqslant 7$ this function is CCZ-inequivalent to the Gold functions, and in the case $n=7$ it is CCZ-inequivalent to any power mapping (and, therefore, to any APN function belonging to one of the families of APN functions known so far).
\end{abstract}

© 2008 Elsevier Inc. All rights reserved.

\section*{1. Introduction}
For any positive integers $n$ and $\delta$, a function $F$ from the field $\mathbb{F}_{2^{n}}$ to itself is called differentially $\delta$-uniform [20] if for every $a \neq 0$ and every $b$ in $\mathbb{F}_{2^{n}}$, the equation $F(x)+F(x+a)=b$ admits at most $\delta$ solutions. Vectorial Boolean functions used as S-boxes in block ciphers must have low differential uniformity to allow high resistance to differential cryptanalysis (see [1]). In this sense, differentially 2-uniform functions, called almost perfect nonlinear (APN), are optimal (since for any function, we have $\delta \geqslant 2$ ).

The notion of APN function is closely connected to the notion of almost bent (AB) function [12] which can be described in terms of the Walsh transform of functions, that we define now in the framework of the field $\mathbb{F}_{2^{n}}$ (it can equivalently be defined in the framework of the vector space $\mathbb{F}_{2}^{n}$, see

\footnotetext{\begin{itemize}
  \item Corresponding author.
\end{itemize}

E-mail addresses: \href{mailto:lilia.b@mail.ru}{lilia.b@mail.ru} (L. Budaghyan), \href{mailto:claude.carlet@inria.fr}{claude.carlet@inria.fr} (C. Carlet), \href{mailto:g.leander@mat.dtu.dk}{g.leander@mat.dtu.dk} (G. Leander).
}
e.g. [8]). Let $\operatorname{tr}(x)$ denote the trace function from $\mathbb{F}_{2^{n}}$ into $\mathbb{F}_{2}$ (that is, $\operatorname{tr}(x)=x+x^{2}+x^{4}+\cdots+x^{2^{n-1}}$ ), then the Walsh transform of a given function $F: \mathbb{F}_{2^{n}} \rightarrow \mathbb{F}_{2^{n}}$ is the integer-valued function
$$
\lambda_{F}(a, b)=\sum_{x \in \mathbb{F}_{2^{n}}}(-1)^{\operatorname{tr}(b F(x)+a x)}, \quad a, b \in \mathbb{F}_{2^{n}} .
$$

The set $\Lambda_{F}=\left\{\lambda_{F}(a, b): a, b \in \mathbb{F}_{2^{n}}, b \neq 0\right\}$ is called the Walsh spectrum of $F$ and the set $\left\{\left|\lambda_{F}(a, b)\right|\right.$ : $\left.a, b \in \mathbb{F}_{2^{n}}, b \neq 0\right\}$ the extended Walsh spectrum of $F$. If the Walsh spectrum of $F$ equals $\left\{0, \pm 2^{\frac{n+1}{2}}\right\}$ then the function $F$ is called $\mathrm{AB}[12]$. AB functions exist for $n$ odd only and oppose an optimum resistance to linear cryptanalysis (see [19]). Besides, every AB function is APN [12], and in the $n$ odd case, any quadratic function is APN if and only if it is $\mathrm{AB}[11]$. A comprehensive survey on APN and AB functions can be found in [10].

The APN and AB properties are preserved by some transformations of functions [11,20]. If $F$ is an APN (respectively AB) function, $A_{1}, A_{2}$ are affine permutations and $A$ is affine then the function $F^{\prime}=A_{1} \circ F \circ A_{2}+A$ is also APN (respectively AB). The functions $F$ and $F^{\prime}$ are then called extended affine equivalent (EA-equivalent). In [11], Carlet, Charpin and Zinoviev introduced an equivalence relation of functions, more recently called CCZ-equivalence, which preserves APN and AB properties: two mappings $F$ and $F^{\prime}$ from $\mathbb{F}_{2^{n}}$ to itself are called CCZ-equivalent if for some affine permutation $\mathcal{L}$ of $\mathbb{F}_{2^{n}}^{2}$ the image of the graph of $F$ is the graph of $F^{\prime}$, that is, $\mathcal{L}\left(G_{F}\right)=G_{F^{\prime}}$ where $G_{F}=\left\{(x, F(x)) \mid x \in \mathbb{F}_{2^{n}}\right\}$ and $G_{F^{\prime}}=\left\{\left(x, F^{\prime}(x)\right) \mid x \in \mathbb{F}_{2^{n}}\right\}$. CCZ-equivalent functions have the same differential uniformity and the same extended Walsh spectrum. EA-equivalence is a particular case of CCZ-equivalence and any permutation is CCZ-equivalent to its inverse. In [8], it is proven that CCZ-equivalence is more general.

There are six known classes of power APN functions, four of which are also AB when $n$ is odd (see [16]). It is conjectured by Dobbertin that any power APN function belongs to one of these six classes [15].

Recently a few classes of quadratic APN polynomials CCZ-inequivalent to power functions have been constructed in [2,3,6,7]. Every of these families is defined under a certain condition on $n$ (such as divisibility of $n$ by a certain number) and none of these families is described by polynomials with coefficients in $\mathbb{F}_{2}$. A natural question is whether all APN polynomials with coefficients in $\mathbb{F}_{2}$ are CCZequivalent to power functions. In the present paper we show that the answer to this question is negative. We give a new approach for constructing quadratic APN functions and using it we construct a class of quadratic APN polynomials with coefficients in $\mathbb{F}_{2}$. We prove that the function $F(x)= x^{3}+\operatorname{tr}\left(x^{9}\right)$ is APN over $\mathbb{F}_{2^{n}}$ for any $n$, and that for all $n \geqslant 7$ it is CCZ-inequivalent to the Gold functions (and, hence, EA-inequivalent to all power functions since Gold functions are the only quadratic APN power functions), to the inverse and Dobbertin functions. Obviously, this function is AB for all odd $n$. We conjecture that for $n \geqslant 7$ the function $F$ is CCZ-inequivalent to any power function. This conjecture is confirmed for the case $n=7$. Further we show that applying CCZ-equivalence to quadratic APN functions, it is possible to construct classes of nonquadratic APN mappings CCZ-inequivalent to power functions.

\section*{2. Construction of new quadratic APN functions}
In the theorem below we give a general approach for constructing new quadratic APN functions from known ones.

Theorem 1. Let $F$ be a quadratic APN function from $\mathbb{F}_{2^{n}}$ to itself, let $f$ be a quadratic Boolean function on $\mathbb{F}_{2^{n}}$ and

$$
\begin{aligned}
\varphi_{F}(x, a) & =F(x)+F(x+a)+F(a)+F(0) \\
\varphi_{f}(x, a) & =f(x)+f(x+a)+f(a)+f(0)
\end{aligned}
$$

If for every nonzero $a \in \mathbb{F}_{2^{n}}$ there exists a linear Boolean function $\ell_{a}$ satisfying the conditions\\
(1) $\varphi_{f}(x, a)=\ell_{a}\left(\varphi_{F}(x, a)\right)$,\\
(2) if $\varphi_{F}(x, a)=1$ for some $x \in \mathbb{F}_{2^{n}}$ then $\ell_{a}(1)=0$,\\
then the function $F(x)+f(x)$ is APN.

Proof. Since the function $F(x)+f(x)$ is quadratic, it is APN if and only if, for every nonzero $a \in \mathbb{F}_{2^{n}}$, the equation $\varphi_{F}(x, a)+\varphi_{f}(x, a)=0$ admits at most two solutions (see e.g. [10]). According to the hypothesis on $\ell_{a}$, a solution to this equation must be such that $\varphi_{f}(x, a)=0$ and therefore such that $\varphi_{F}(x, a)=0$. Then, $F$ being quadratic APN, this equation admits at most two solutions.

The same principle as in Theorem 1 allows generating a large variety of differentially 4 -uniform functions from APN functions as it is shown in the proposition below.

Proposition 1. For any APN function F the following functions are differentially 4-uniform\\
(1) $F(x)+\operatorname{tr}(G(x))$ for any function $G$;\\
(2) $F \circ A$ and $A \circ F$ for any affine function $A$ which is 2-to-1;\\
(3) a function $F^{\prime}$ such that $\mathcal{L}\left(G_{F}\right)=G_{F^{\prime}}$ for some affine 2-to-1 mapping $\mathcal{L}$ of $\mathbb{F}_{2^{n}}^{2}$.

Remark 1. Note that, for every APN quadratic function $F$, every quadratic Boolean function $f$ and every nonzero vector $a$, a linear function $l_{a}$ satisfying $\varphi_{f}(x, a)=\ell_{a}\left(\varphi_{F}(x, a)\right)$ always exists. This is due to the fact that, by the assumption that $F$ is APN, the kernel of $\varphi_{F}(x, a)$ equals $\{0, a\}$. This set is always a subset of the kernel of $\varphi_{f}(x, a)$, and this is the necessary and sufficient condition for the existence of $l_{a}$.

A direct consequence of Theorem 1 is that, if $F$ is APN and if $\ell$ is a linear form such that $\ell(1)=0$, then the function $F(x)+\ell(F(x))$ is APN. But this function is affine equivalent to $F$ since it is equal to $L \circ F$ where $L(x)=x+\ell(x)$, and the condition that $\ell(1)=0$ is equivalent to saying that $L$ is a permutation.

We give now an example where Theorem 1 leads to a function which is CCZ-inequivalent to the original function $F$.

Corollary 1. Let $n$ be any positive integer. Then the function $x^{3}+\operatorname{tr}\left(x^{9}\right)$ is APN on $\mathbb{F}_{2^{n}}$.\\
Proof. We can apply Theorem 1 with $F(x)=x^{3}, \varphi_{F}(x, a)=a^{2} x+a x^{2}, f(x)=\operatorname{tr}\left(x^{9}\right), \varphi_{f}(x, a)=\operatorname{tr}\left(a^{8} x+\right. \left.a x^{8}\right)$ and $\ell_{a}(y)=\operatorname{tr}\left(a^{6} y+a^{3} y^{2}+a^{-3} y^{4}\right)$. Indeed, we have then

$$
\ell_{a}\left(\varphi_{F}(x, a)\right)=\operatorname{tr}\left(a^{6}\left(a^{2} x+a x^{2}\right)+a^{3}\left(a^{4} x^{2}+a^{2} x^{4}\right)+a^{-3}\left(a^{8} x^{4}+a^{4} x^{8}\right)\right)=\varphi_{f}(x, a)
$$

and if there exists $x \in \mathbb{F}_{2}^{n}$ such that $\varphi_{F}(x, a)=1$ then

$$
\ell_{a}(1)=\operatorname{tr}\left(a^{-3}\right)=\operatorname{tr}\left(\frac{x}{a}+\left(\frac{x}{a}\right)^{2}\right)=0 .
$$

Remark 2. The APN property of the function $x^{3}+\operatorname{tr}\left(x^{9}\right)$ can be proven also with the following arguments due to Dillon [14]. If $F$ is a quadratic function then for any nonzero $a$ and for $\varphi_{F}(x, y)= F(x+y)+F(x)+F(y)+F(0)$ there exists a linear function $L_{a}$ such that $\varphi_{F}(a x, a)=L_{a}\left(x+x^{2}\right)$. Indeed, if $F(x)=\sum_{i \leqslant j} c_{i, j} x^{2^{i}+2^{j}}$ then $L_{a}(z)=\sum_{i \leqslant j} c_{i, j} a^{2^{i}+2^{j}}\left(T_{j-i}(z)\right)^{2^{i}}$, where $T_{k}(z)=z+z^{2}+\cdots+z^{2^{k-1}}$. Thus, $F$ is APN if and only if for any nonzero $a$ and $z$ the equality $L_{a}(z)=0$ implies $\operatorname{tr}(z)=1$. In the case when $F(x)=x^{3}+\operatorname{tr}\left(x^{9}\right)$ we have $L_{a}(z)=a^{3} z+\operatorname{tr}\left(a^{9} T_{3}(z)\right)$ and if $L_{a}(z)=0$ for some $z \neq 0$ then $1=a^{3} z=\operatorname{tr}\left(a^{9} T_{3}(z)\right)$ which implies $1=\operatorname{tr}\left(z^{-3}\left(z+z^{2}+z^{4}\right)\right)=\operatorname{tr}(z)$.

Another class of APN functions $F$, to which the construction of Theorem 1 can be applied, is a class of trinomial APN functions described in [6]. However, for this class of functions we were able to construct only functions that are EA-equivalent to the original trinomial $F$. More precisely we have the following proposition.

Proposition 2. Let $m$ be a positive odd integer, $n=2 m$, $\alpha$ a primitive element of $\mathbb{F}_{2^{n}}$. Then the functions $F(x)=x^{6}+x^{2^{m}+1}+\alpha^{2^{m}-1} x^{6 \cdot 2^{m}}$ and $G(x)=F(x)+\operatorname{tr}\left(\alpha^{2^{m-1}+1} x^{3}\right)$ are EA-equivalent on $\mathbb{F}_{2^{n}}$.

Proof. Let $u=\frac{\alpha^{2^{m+1}+1}}{\alpha^{2^{m}-1}+1}$ and $L(x)=x+\operatorname{tr}(u x)$. It is not difficult to see that $L(F(x))=G(x)$.\\
There is a straightforward generalization of Theorem 1:\\
Theorem 2. Let $F$ be a quadratic APN function from $\mathbb{F}_{2^{n}}$ to itself, let $f$ be a quadratic function from $\mathbb{F}_{2^{n}}$ to $\mathbb{F}_{2^{m}}$ where $m$ is a divisor of $n$, and

$$
\begin{aligned}
\varphi_{F}(x, a) & =F(x)+F(x+a)+F(a)+F(0) \\
\varphi_{f}(x, a) & =f(x)+f(x+a)+f(a)+f(0)
\end{aligned}
$$

If for every nonzero $a \in \mathbb{F}_{2^{n}}$ there exists a linear function $\ell_{a}$ from $\mathbb{F}_{2^{n}}$ to $\mathbb{F}_{2^{m}}$ which satisfies the conditions\\
(1) $\varphi_{f}(x, a)=\ell_{a}\left(\varphi_{F}(x, a)\right)$,\\
(2) for every $u \in \mathbb{F}_{2^{m}}^{*}$, if $\varphi_{F}(x, a)=u$ for some $x \in \mathbb{F}_{2^{n}}$ then $\ell_{a}(u) \neq u$,\\
then the function $F(x)+f(x)$ is APN.\\
We could find an application of Theorem 2:\\
Corollary 2. Let $n=2 m$ where $m$ is an even positive integer. Let us denote by $\operatorname{tr}_{n / m}$ the trace function from $\mathbb{F}_{2^{n}}$ to $\mathbb{F}_{2^{m}}: \operatorname{tr}_{n / m}(x)=x+x^{2^{m}}$. The functions $F(x)=x^{3}+\operatorname{tr}_{n / m}\left(x^{2^{m}+2}\right)=x^{3}+x^{2^{m}+2}+x^{2^{m+1}+1}$ and $F^{\prime}(x)= x^{3}+\left(\operatorname{tr}_{n / m}(x)\right)^{3}$ are APN.

But unfortunately, these functions are EA-equivalent to power functions. Indeed, let $G$ be the Gold function $G(x)=x^{2^{m-1}+1}$. Let $\gamma$ be any element of $\mathbb{F}_{4} \backslash \mathbb{F}_{2}$ and $L_{1}, L_{2}$ be the linear mappings $L_{1}(x)= \gamma^{2} x^{2^{m+1}}+\gamma x^{2}, L_{2}(x)=\gamma x^{2^{m}}+\gamma^{2} x$. Then $G \circ L_{1}(x)=L_{2} \circ F(x)$. Note further that $F^{\prime}(x)=(F(x))^{2^{m}}$.

\subsection*{2.1. An algorithmic approach}
Below we describe an algorithmic approach to search for functions fulfilling the conditions of Theorem 1 when $F$ is a Gold function. The first step will be to find an explicit description of the linear function $l_{a}$ used in Theorem 1. Let $F(x)=x^{2^{r}+1}$ and $f(x)=\operatorname{tr}\left(\alpha x^{2^{i}+1}\right)$. Then

$$
\begin{gathered}
\varphi_{F}(x, a)=a^{2^{r}+1}\left(\left(\frac{x}{a}\right)+\left(\frac{x}{a}\right)^{2^{r}}\right), \\
\varphi_{f}(x, a)=\operatorname{tr}\left(a^{2^{i}+1} \alpha\left(\left(\frac{x}{a}\right)+\left(\frac{x}{a}\right)^{2^{i}}\right)\right) .
\end{gathered}
$$

If we define $t=\left(i r^{-1}-1\right) \bmod n$ we get

$$
\varphi_{f}(x, a)=\operatorname{tr}\left(a^{2^{i}+1} \alpha\left(\sum_{j=0}^{t}\left[\frac{\varphi_{F}(x, a)}{a^{2^{r}+1}}\right]^{2^{j r}}\right)\right)=\operatorname{tr}\left(\left(\sum_{j=0}^{t} \alpha^{2^{-j r}} a^{2^{i-j r}+2^{-j r}-\left(2^{r}+1\right)}\right) \varphi_{F}(x, a)\right) .
$$

Thus denoting $T_{i}^{r}(a, \alpha)=\sum_{j=0}^{t} \alpha^{2^{-j r}} a^{2^{i-j r}+2^{-j r}-\left(2^{r}+1\right)}$ we get

$$
\varphi_{f}(x, a)=\operatorname{tr}\left(T_{i}^{r}(a, \alpha) \varphi_{F}(x, a)\right)
$$

In general for $g(x)=\sum_{i} \alpha_{i} x^{2^{i}+1}$ we get

$$
\varphi_{g}(x, a)=\operatorname{tr}\left(\left(\sum_{i} T_{i}^{r}\left(a, \alpha_{i}\right)\right) \varphi_{F}(x, a)\right) .
$$

Following Theorem 1, the condition for $F+g$ to be APN is that, if $\operatorname{tr}\left(a^{-\left(2^{r}+1\right)}\right)=0$ then

$$
\operatorname{tr}\left(\sum_{i} T_{i}^{r}\left(a, \alpha_{i}\right)\right)=\sum_{i} \operatorname{tr}\left(T_{i}^{r}\left(a, \alpha_{i}\right)\right)=0 .
$$

Fixing a base $\left(b_{j}\right)_{j=1, \ldots, n}$ of $\mathbb{F}_{2^{n}}$ over $\mathbb{F}_{2}$ we can consider the set of vectors

$$
\left\{\operatorname{tr}\left(T_{i}^{r}\left(a, b_{j}\right)\right)_{a \in \mathbb{F}_{2^{n}}, \operatorname{tr}\left(a^{-3}\right)=0} \mid i, j \in\{0 \ldots n-1\}\right\}
$$

Given $F$, finding a quadratic function $g$ such that the conditions of Theorem 1 are fulfilled is equivalent to finding a set of linearly dependent vectors in this set. We computed these vectors and all linear dependent sets up to dimension 15. The only examples in addition to $x^{3}+\operatorname{tr}\left(x^{9}\right)$ are listed below.\\
(1) For $n=5$ the function $x^{5}+\operatorname{tr}\left(x^{3}\right)$ is APN (CCZ-equivalent to Gold functions).\\
(2) For $n=8$ the function $x^{9}+\operatorname{tr}\left(x^{3}\right)$ is APN (CCZ-inequivalent to power functions and to $x^{3}+\operatorname{tr}\left(x^{9}\right)$ ).

\section*{3. CCZ-inequivalence of the new APN function to other known APN mappings}
Any of the known families of APN functions except the family of the function $x^{3}$ cannot cover the family of the function $x^{3}+\operatorname{tr}\left(x^{9}\right)$. Indeed, by Corollary 1 the function $x^{3}+\operatorname{tr}\left(x^{9}\right)$ is APN over $\mathbb{F}_{2^{n}}$ for any $n$ while all previously known APN functions (power mappings as well as polynomials) are defined under some condition on $n$, except the function $x^{3}$ which is APN for any $n$. Hence, when proving CCZinequivalence of $x^{3}+\operatorname{tr}\left(x^{9}\right)$ to the known APN functions, we mainly concentrate our attention on the Gold functions. Below we prove that $x^{3}+\operatorname{tr}\left(x^{9}\right)$ is CCZ-inequivalent to Gold mappings for any $n \geqslant 7$. Besides, we prove that for $n=7$ the function $x^{3}+\operatorname{tr}\left(x^{9}\right)$ is CCZ-inequivalent to any function from the known classes of APN mappings.

Theorem 3. The function $x^{3}+\operatorname{tr}\left(x^{9}\right)$ is CCZ-inequivalent to any Gold function on $\mathbb{F}_{2^{n}}$ when $n \geqslant 7$.\\
Proof. Let $F(x)=x^{3}+\operatorname{tr}\left(x^{9}\right)$ and $G(x)=x^{2^{r}+1}$ over $\mathbb{F}_{2^{n}}$ with $n \geqslant 7, r \leqslant(n-1) / 2, \operatorname{gcd}(r, n)=1$. We also assume further that $n \neq 8,10$. For $n=8,10$ CCZ-inequivalence of $F$ and $G$ is checked by a computer.

Suppose first that the functions $F$ and $G$ are EA-equivalent. Then, there exist affine permutations $L_{1}, L_{2}$ and an affine function $L^{\prime}$ such that

$$
L_{1}\left(x^{3}\right)+L_{1}\left(\operatorname{tr}\left(x^{9}\right)\right)=\left(L_{2}(x)\right)^{2^{t}+1}+L^{\prime}(x) .
$$

That is,

$$
L_{1}\left(x^{3}\right)+L_{1}(1) \operatorname{tr}\left(x^{9}\right)=\left(L_{2}(x)\right)^{2^{t}+1}+L^{\prime}(x)
$$

Since the functions are quadratic, we can assume without loss of generality that $L_{1}$ and $L_{2}$ are linear: $L_{1}(x)=\sum_{m \in \mathbb{Z} / n \mathbb{Z}} b_{m} x^{2^{m}}, L_{2}(x)=\sum_{p \in \mathbb{Z} / n \mathbb{Z}} c_{p} x^{2^{p}}$. Then we get


\begin{equation*}
\sum_{m \in \mathbb{Z} / n \mathbb{Z}} b_{m} x^{3 \cdot 2^{m}}+\operatorname{tr}\left(x^{9}\right) \sum_{m \in \mathbb{Z} / n \mathbb{Z}} b_{m}=\sum_{l, p \in \mathbb{Z} / n \mathbb{Z}} c_{p} c_{l}^{2^{t}} x^{2^{l+t}+2^{p}}+L^{\prime}(x) \tag{1}
\end{equation*}


On the left-hand side of the identity (1) we have only items of the type $x^{3 \cdot 2^{m}}, x^{9 \cdot 2^{m}}$, with some coefficients. Therefore this must be true also for the right-hand side of the identity. Since $n \geqslant 7$ and $n \neq 8,10$ then there exists a positive integer $p$ different from 1 and 3 such that $\operatorname{gcd}(n, p)=1$ and $n>2 p$. Since $n>2 p$ then $2^{p}+1$ is not in the same cyclotomic coset with 3 or 9 . Therefore, the items of the type $x^{2^{k}\left(2^{p}+1\right)}$ must cancel. That is, for any $k$


\begin{equation*}
c_{k} c_{k-t+p}^{2^{t}}=c_{k+p} c_{k-t}^{2^{t}} \tag{2}
\end{equation*}


Since $n \geqslant 7$ then 3 and 9 are in different cyclotomic cosets and we have for any $k$

$$
L_{1}(1)=c_{k} c_{k-t+3}^{2^{t}}+c_{k+3} c_{k-t}^{2^{t}}
$$

If $L_{1}(1) \neq 0$ then


\begin{equation*}
c_{k} c_{k-t+3}^{2^{t}} \neq c_{k+3} c_{k-t}^{2^{t}} \tag{3}
\end{equation*}


If $c_{k} \neq 0$ for all $k$ then from (2) and (3) we get


\begin{align*}
& c_{k} c_{k-t}^{-2^{t}}=c_{k+p} c_{k-t+p}^{-2^{t}}  \tag{4}\\
& c_{k} c_{k-t}^{-2^{t}} \neq c_{k+3} c_{k-t+3}^{-2^{t}} \tag{5}
\end{align*}


Since $\operatorname{gcd}(n, p)=1$ and from (4) we obtain $c_{k} c_{k-t}^{-2^{t}}=c_{m} c_{m-t}^{-2^{t}}$ for any $m$. It contradicts (5). Thus, $c_{k}=0$ for some $k$. Then from (2) and (3) we get that $c_{k+p}=0$. Repeating this step for $c_{k+p}, c_{k+2 p}, \ldots$ we get $c_{k+p s}=0$ and since $\operatorname{gcd}(n, p)=1$ then $c_{k}=0$ for all $k$. A contradiction. If $L_{1}(1)=0$ then the equation $L(x)=0$ has at least 2 solutions 0,1 and therefore $L_{1}$ is not a permutation. Thus, $F$ and $G$ are EA-inequivalent.

Suppose that $F(x)$ and $G(x)$ are CCZ-equivalent. Then there exists an affine automorphism $L= \left(L_{1}, L_{2}\right)$ of $\mathbb{F}_{2^{n}} \times \mathbb{F}_{2^{n}}$ such that

$$
y=F(x) \quad \Leftrightarrow \quad L_{2}(x, y)=G\left(L_{1}(x, y)\right)
$$

and $L_{1}(x, F(x))$ is a permutation (see [7]). This implies then

$$
L_{2}(x, F(x))=G\left(L_{1}(x, F(x))\right)
$$

Writing $L_{1}(x, y)=L(x)+L^{\prime}(y)$ and $L_{2}(x, y)=L^{\prime \prime}(x)+L^{\prime \prime \prime}(y)$ gives


\begin{equation*}
L^{\prime \prime}(x)+L^{\prime \prime \prime}(F(x))=G\left(L(x)+L^{\prime}(F(x))\right) \tag{6}
\end{equation*}


We can write

$$
\begin{aligned}
& L(x)=b+\sum_{m \in \mathbb{Z} / n \mathbb{Z}} b_{m} x^{2^{m}} \\
& L^{\prime}(x)=b^{\prime}+\sum_{m \in \mathbb{Z} / n \mathbb{Z}} b_{m}^{\prime} x^{2^{m}} \\
& L^{\prime \prime}(x)=b^{\prime \prime}+\sum_{m \in \mathbb{Z} / n \mathbb{Z}} b_{m}^{\prime \prime} x^{2^{m}} \\
& L^{\prime \prime \prime}(x)=b^{\prime \prime \prime}+\sum_{m \in \mathbb{Z} / n \mathbb{Z}} b_{m}^{\prime \prime \prime} x^{2^{m}} \\
& b+b^{\prime}=c
\end{aligned}
$$

Then we get

$$
\begin{aligned}
G\left(L(x)+L^{\prime}(F(x))\right)= & \left(L(x)+L^{\prime}\left(x^{3}+\operatorname{tr}\left(x^{9}\right)\right)\right)\left(L(x)+L^{\prime}\left(x^{3}+\operatorname{tr}\left(x^{9}\right)\right)\right)^{2^{r}} \\
= & {\left[\sum_{m, k \in \mathbb{Z} / n \mathbb{Z}} b_{m} b_{k}^{2^{r}} x^{2^{m}}+2^{k+r}+2^{k+r+1}+L^{\prime}(1)^{2^{r}} \sum_{m, k \in \mathbb{Z} / n \mathbb{Z}} b_{m} x^{2^{m}+2^{k+3}+2^{k}}\right.} \\
& \left.+\sum_{m, k \in \mathbb{Z} / n \mathbb{Z}} b_{m}^{\prime} b_{k}^{2^{r}} x^{2^{m+1}+2^{m}+2^{k+r}}+L^{\prime}(1) \sum_{m, k \in \mathbb{Z} / n \mathbb{Z}} b_{m}^{2^{r}} x^{2^{m+r}+2^{k+3}+2^{k}}\right] \\
& +\left[\sum_{m, k \in \mathbb{Z} / n \mathbb{Z}} b_{m}^{\prime} b_{k}^{\prime 2^{r}} x^{2^{m+1}+2^{m}+2^{k+r+1}+2^{k+r}}+L^{\prime}(1)^{2^{r}} \sum_{m, k \in \mathbb{Z} / n \mathbb{Z}} b_{m}^{\prime} x^{2^{m+1}+2^{m}+2^{k+3}+2^{k}}\right. \\
& \left.+L^{\prime}(1) \sum_{m, k \in \mathbb{Z} / n \mathbb{Z}} b_{m}^{\prime 2^{r}} x^{2^{m+r+1}+2^{m+r}+2^{k+3}+2^{k}}\right]+Q(x)
\end{aligned}
$$

where $Q(x)$ is a quadratic polynomial. Obviously, all terms in the expression above whose exponents have 2 -weight strictly greater than 2 must cancel.

Since $F$ and $G$ are EA-inequivalent then $L^{\prime}$ is not a constant. Then there exists $m \in \mathbb{Z} / n \mathbb{Z}$ such that $b_{m}^{\prime} \neq 0$.

Let $L^{\prime}(1) \neq 0$. Since the items with the exponent $2^{m+1}+2^{m}+2^{m+2}+2^{m+5}$ have to vanish then we get $L^{\prime}(1)^{2^{r}} b_{m}^{\prime}=L^{\prime}(1) b_{m-r}^{2^{r}}$ and since $L^{\prime}(1) \neq 0, b_{m}^{\prime} \neq 0$ and $r$ is coprime with $n$ then $b_{k}^{\prime} \neq 0$ and $b_{k}^{\prime} b_{k-r}^{\prime-2^{r}}=L^{\prime}(1)^{1-2^{r}}$ for all $k$. Now we can deduce that $b_{k+r}^{\prime}=L^{\prime}(1)^{1-2^{r}} b_{k}^{\prime 2^{r}}$ for all $k$. Then, introducing $\mu$ such that $L^{\prime}(1)^{1-2^{r}}=\mu^{2^{r}-1}$, we deduce that $\mu b_{k+r}^{\prime}=\left(\mu b_{k}^{\prime}\right)^{2^{r}}$ for all $k$ and then that $\mu b_{k+1}^{\prime}= \left(\mu b_{k}^{\prime}\right)^{2}$ (using that $\operatorname{gcd}(r, n)=1$ ) and then $\mu b_{k}^{\prime}=\left(\mu b_{0}^{\prime}\right)^{2^{k}}$. This means that $\mu L^{\prime}(x)=\mu b^{\prime}+\operatorname{tr}\left(\mu b^{\prime}{ }_{0} x\right)$. It implies that all nonquadratic items in the last bracket vanish and $L^{\prime}(x)=d+\operatorname{tr}\left(d^{\prime} x\right)$ for some $d, d^{\prime}$. Thus, the function $L$ is not 0 since $L^{\prime}$ is not a permutation, and then $b_{m} \neq 0$ for some $m$. Since the items with the exponent $2^{m}+2^{m+2}+2^{m+5}$ have to vanish then $L^{\prime}(1)^{2^{r}} b_{m}=L^{\prime}(1) b_{m-r}^{2^{r}}$. Like above we get $L(x)=d+\operatorname{tr}\left(d^{\prime} x\right)$. Thus, $L_{1}(x, F(x))=d^{\prime \prime}+\operatorname{tr}\left(F^{\prime}(x)\right)$ for some $d^{\prime \prime}$ and $F^{\prime}(x)$ and $L_{1}(x, F(x))$ is not a permutation. A contradiction.

Let $L^{\prime}(1)=0$ and $r \neq 1$. Then $2^{m+1}+2^{m}+2^{m+r+1}+2^{m+r}$ has 2 -weight 4 and since the items with this exponent should cancel then we get ${b_{m}^{\prime}}^{2^{r}+1}=b_{m+r}^{\prime}{b_{m-r}^{\prime}}^{2^{r}}$. Since $b_{m}^{\prime} \neq 0$ then $b_{m+r}^{\prime}, b_{m-r}^{\prime} \neq 0$ and $b_{m}^{\prime} b_{m-r}^{\prime-2^{r}}=b_{m+r}^{\prime} b_{m}^{\prime-2^{r}}$. Since $\operatorname{gcd}(n, r)=1$ then $b_{k}^{\prime} \neq 0, b_{k}^{\prime} b_{k-r}^{\prime-2^{r}}=b_{m}^{\prime} b_{m-r}^{\prime-2^{r}}$ for all $k$ and this implies $L^{\prime}(x)=d+\operatorname{tr}\left(d^{\prime} x\right)$ for some $d, d^{\prime}$. Since $L_{1}(x, F(x))$ is a permutation then $L \neq 0$ and $b_{m} \neq 0$ for some $m$. The items with the exponent $2^{m}+2^{m+r}+2^{m+r+1}$ should vanish. Therefore, $b_{m} b_{m}^{2^{r}}=b_{m+r}^{\prime} b_{m-r}^{2^{r}}$ and\\
$b_{m} b_{m-r}^{-2^{r}}=b_{m+r}^{\prime} b_{m}^{\prime-2^{r}}$. As above it leads to the equality $L(x)=d+\operatorname{tr}\left(d^{\prime} x\right)$ which is in contradiction with $L_{1}(x, F(x))$ being a permutation.

Let $L^{\prime}(1)=0$ and $r=1$. Since $L^{\prime}(1)=0$ and $b_{m}^{\prime} \neq 0$ then there exists $t$ such that $b_{m+t}^{\prime} \neq 0$. If $t \neq-1,-2$ then $2^{m+1}+2^{m}+2^{m+t+2}+2^{m+t+1}$ has 2 -weight 4 and we get $b_{m}^{\prime} b_{m+t}^{\prime 2^{r}}=b_{m+t+1}^{\prime} b_{m-1}^{\prime 2^{r}}$ and $b_{m}^{\prime} b_{m-1}^{\prime-2^{r}}=b_{m+t+1}^{\prime} b_{m+t}^{\prime-2^{r}}$. Therefore, $L^{\prime}(x)=d+\operatorname{tr}\left(d^{\prime} x\right)$ for some $d, d^{\prime}$. If $t \neq 1,2$ then $2^{m+t+1}+2^{m+t}+ 2^{m+2}+2^{m+1}$ has 2 -weight 4 and we get $b_{m+t}^{\prime} b_{m}^{2^{r}}=b_{m+1}^{\prime} b_{m+t-1}^{\prime 2^{r}}$ and again $L^{\prime}(x)=d+\operatorname{tr}\left(d^{\prime} x\right)$ for some $d, d^{\prime}$. Thus, $L \neq 0$ and then $b_{m} \neq 0$ for some $m$. Since the items with the exponent $2^{m}+2^{m+2}+2^{m+3}$ cancel then $b_{m} b_{m+1}^{\prime 2^{r}}=b_{m+2}^{\prime} b_{m-1}^{2^{r}}$ and $b_{m} b_{m-1}^{-2^{r}}=b_{m+2}^{\prime} b_{m+1}^{\prime-2^{r}}$. This implies $L(x)=d+\operatorname{tr}\left(d^{\prime} x\right)$ and, thus, $L_{1}(x, F(x))$ is not a permutation. Therefore, $F$ and $G$ are not CCZ-equivalent.

Corollary 3. The function $x^{3}+\operatorname{tr}\left(x^{9}\right)$ is EA-inequivalent to any power function on $\mathbb{F}_{2^{n}}$ when $n \geqslant 7$.

Proof. The function $x^{3}+\operatorname{tr}\left(x^{9}\right)$ is quadratic APN and by Theorem 3 it is EA-inequivalent to any quadratic power function. Since the algebraic degree is EA-invariant then $F$ is EA-inequivalent to any power mapping.

Dobbertin and inverse APN functions have unique Walsh spectra (except for the case $n=3$ where the inverse function is EA-equivalent to $x^{3}$ ) which are different from the Walsh spectra of quadratic APN functions (see $[9,11,18,21]$ ). Since the extended Walsh spectrum of a function is invariant under CCZ-equivalence, then the function of Corollary 1 is CCZ-inequivalent to the inverse and Dobbertin APN functions for $n \geqslant 7$.

Let $G=\mathbb{F}_{2}\left[\mathbb{F}_{2^{n}} \times \mathbb{F}_{2^{n}}\right]$ be the group algebra of $\mathbb{F}_{2^{n}} \times \mathbb{F}_{2^{n}}$ over $\mathbb{F}_{2}$. It consists of the formal sums $\sum_{g \in G} a_{g} g$ where $a_{g} \in \mathbb{F}_{2}$. If $S$ is a subset of $\mathbb{F}_{2^{n}} \times \mathbb{F}_{2^{n}}$ then it can be identified with the element $\sum_{s \in S} S$ of $G$. For any APN mapping $F$ we denote

$$
\Delta_{F}=\{(a, b): F(x)+F(x+a)=b \text { has } 2 \text { solutions }\} \subset \mathbb{F}_{2^{n}} \times \mathbb{F}_{2^{n}} .
$$

The dimensions of the ideals of $G$ generated by $\Delta_{F}$ and by the graph $G_{F}$ of $F$ are called $\Delta$ - and $\Gamma$-ranks, respectively. According to $[4,13], \Delta$ - and $\Gamma$-ranks of a function are CCZ-invariant.

For $n=7$ the $\Delta$-rank of the function $F(x)=x^{3}+\operatorname{tr}\left(x^{9}\right)$ equals 212 and differs from the $\Delta$-ranks of the Kasami functions $x^{13}$ and $x^{23}$ (which equal 338 and 436, respectively). Thus, for $n=7$ the function $F$ is CCZ-inequivalent to Kasami functions, and by Theorem 3 to the Gold functions. Since in this field the Welch and Niho cases coincide with the Kasami cases then $F$ is CCZ-inequivalent to all power maps on $\mathbb{F}_{2^{7}}$. Besides, all previously known families of APN polynomials are not defined for $n=7$ (see [2,6,7]). Hence we can make the following claim.

Corollary 4. On $\mathbb{F}_{2^{7}}$ the function $F(x)=x^{3}+\operatorname{tr}\left(x^{9}\right)$ is CCZ-inequivalent to power functions and to any known APN mapping.

Applying CCZ-equivalence to the quadratic APN function $F(x)=x^{3}+\operatorname{tr}\left(x^{9}\right)$, it is possible to construct classes of nonquadratic APN mappings which are CCZ-inequivalent to power functions.

For $n$ divisible by 3 we denote $\operatorname{tr}_{n / 3}(x)=x+x^{2^{3}}+x^{2^{6}}+\cdots+x^{2^{(n / 3-1) 3}}$ the trace function of $\mathbb{F}_{2^{n}}$ into $\mathbb{F}_{2^{3}}$.

Proposition 3. Let $F: \mathbb{F}_{2^{n}} \rightarrow \mathbb{F}_{2^{n}}, F(x)=x^{3}+\operatorname{tr}\left(x^{9}\right)$ then the following functions are CCZ-equivalent to $F$\\
(1) for $n$ odd the function with algebraic degree 3

$$
x^{3}+\operatorname{tr}\left(x^{9}\right)+\left(x^{2}+x\right) \operatorname{tr}\left(x^{3}+x^{9}\right)
$$

(2) for $n$ even the function with algebraic degree 3

$$
x^{3}+\operatorname{tr}\left(x^{9}\right)+\left(x^{2}+x+1\right) \operatorname{tr}\left(x^{3}\right)
$$

(3) for $n$ divisible by 6 the function with algebraic degree 4

$$
\left(x+\operatorname{tr}_{n / 3}\left(x^{6}+x^{12}\right)+\operatorname{tr}(x) \operatorname{tr}_{n / 3}\left(x^{3}+x^{12}\right)\right)^{3}+\operatorname{tr}\left(\left(x+\operatorname{tr}_{n / 3}\left(x^{6}+x^{12}\right)+\operatorname{tr}(x) \operatorname{tr}_{n / 3}\left(x^{3}+x^{12}\right)\right)^{9}\right) ;
$$

(4) for $n$ odd and divisible by 3 the function with algebraic degree 4

$$
\left(x^{\frac{1}{3}}+\operatorname{tr}_{n / 3}\left(x+x^{4}\right)\right)^{-1}+\operatorname{tr}\left(\left(\left(x^{\frac{1}{3}}+\operatorname{tr}_{n / 3}\left(x+x^{4}\right)\right)^{-1}\right)^{3}\right) .
$$

Proof. The proof is the same as for the cases from $[5,8]$ (use the affine permutation $\mathcal{L}(x, y)=(x+ \operatorname{tr}(y), y)$ for the first two cases, $\mathcal{L}(x, y)=\left(x+\operatorname{tr}_{n / 3}\left(y^{2}+y^{4}\right), y\right)$ for the third case and $\mathcal{L}(x, y)= \left(x+\operatorname{tr}_{n / 3}\left(y+y^{4}\right), y\right)$ for the fourth case $)$.

Remark 3. Note that the second and the fourth APN functions in Proposition 3 can be obtained from the functions of $[5,8]$ by adding $\operatorname{tr}(G(x))$ for some $G$. It means that the construction $F(x)+\operatorname{tr}(G(x))$, where $F$ is APN and $G$ is arbitrary, actually gives new APN functions even if $F$ and $G$ are not quadratic.

Note added in proof: After the appearance of the present paper as a preprint, the idea of adding a Boolean function to a known APN function for obtaining a new APN function has been further used in [17] where an APN function over $\mathbb{F}_{2^{6}}$ CCZ-inequivalent to power functions and to quadratic functions is found by computer investigations. Finding a method producing an infinite class of such APN functions (in particular, which includes this isolated function) is an open question.

\section*{Acknowledgments}
We would like to thank John Dillon for nice observations, and Willem De Graaf and Nobuo Nakagawa for useful discussions. We are also grateful to a reviewer of this paper for advice on improving the formulation of Theorem 3. The research of the first author was partly supported by a postdoctoral fellowship of MIUR-Italy via PRIN 2006, also by Lavoro eseguito con il contributo della Provincia autonoma di Trento; also by Norwegian Research Council. The research of the third author was supported by a DAAD postdoctoral fellowship.

\section*{References}
[1] E. Biham, A. Shamir, Differential cryptanalysis of DES-like cryptosystems, J. Cryptology 4 (1) (1991) 3-72.\\[0pt]
[2] C. Bracken, E. Byrne, N. Markin, G. McGuire, New families of quadratic almost perfect nonlinear trinomials and multinomials, Finite Fields Appl. 14 (2008) 703-714.\\[0pt]
[3] C. Bracken, E. Byrne, N. Markin, G. McGuire, A few more quadratic APN functions, preprint, 2008. Available at \href{http://arxiv}{http://arxiv}. org/abs/0804.4799.\\[0pt]
[4] K. Browning, J.F. Dillon, R.E. Kibler, M. McQuistan, APN polynomials and related codes, to appear in a special volume of J. Combin. Inform. System Sci., 2008, in press; honoring the 75th birthday of Prof. D.K. Ray-Chaudhuri.\\[0pt]
[5] L. Budaghyan, The simplest method for constructing APN polynomials EA-inequivalent to power functions, in: Proceedings of WAIFI 2007, in: Lecture Notes in Comput. Sci., vol. 4547, 2007, pp. 177-188.\\[0pt]
[6] L. Budaghyan, C. Carlet, Classes of quadratic APN trinomials and hexanomials and related structures, IEEE Trans. Inform. Theory 54 (5) (May 2008) 2354-2357.\\[0pt]
[7] L. Budaghyan, C. Carlet, G. Leander, Two classes of quadratic APN binomials inequivalent to power functions, IEEE Trans. Inform. Theory 54 (9) (Sept. 2008) 4218-4229.\\[0pt]
[8] L. Budaghyan, C. Carlet, A. Pott, New classes of almost bent and almost perfect nonlinear functions, IEEE Trans. Inform. Theory 52 (3) (March 2006) 1141-1152.\\[0pt]
[9] A. Canteaut, P. Charpin, H. Dobbertin, Weight divisibility of cyclic codes, highly nonlinear functions on $\mathbb{F}_{2^{m}}$, and crosscorrelation of maximum-length sequences, SIAM J. Discrete Math. 13 (1) (2000) 105-138.\\[0pt]
[10] C. Carlet, Vectorial Boolean functions for cryptography, chapter of the monography Boolean Methods and Models, Y. Crama, P. Hammer (Eds.), Cambridge Univ. Press, in press.\\[0pt]
[11] C. Carlet, P. Charpin, V. Zinoviev, Codes, bent functions and permutations suitable for DES-like cryptosystems, Des. Codes Cryptogr. 15 (2) (1998) 125-156.\\[0pt]
[12] F. Chabaud, S. Vaudenay, Links between differential and linear cryptanalysis, in: Advances in Cryptology - EUROCRYPT'94, in: Lecture Notes in Comput. Sci., vol. 950, Springer-Verlag, New York, 1995, pp. 356-365.\\[0pt]
[13] J.F. Dillon, APN Polynomials and Related Codes, Polynomials over Finite Fields and Applications, Banff International Research Station, Nov. 2006.\\[0pt]
[14] J.F. Dillon, private communication, Feb. 2007.\\[0pt]
[15] H. Dobbertin, Almost perfect nonlinear power functions over $G F\left(2^{n}\right)$ : the Niho case, Inform. and Comput. 151 (1999) 57-72.\\[0pt]
[16] H. Dobbertin, Almost perfect nonlinear power functions over $G F\left(2^{n}\right)$ : a new case for $n$ divisible by 5, in: Proceedings of Finite Fields and Applications FQ5, Augsburg, Germany, Springer, 2000, pp. 113-121.\\[0pt]
[17] Y. Edel, A. Pott, A new almost perfect nonlinear function which is not quadratic, preprint, 2008. Available at \href{http://eprint}{http://eprint}. \href{http://iacr.org/2008/313}{iacr.org/2008/313}.\\[0pt]
[18] G. Lachaud, J. Wolfmann, The weights of the orthogonals of the extended quadratic binary Goppa codes, IEEE Trans. Inform. Theory 36 (1990) 686-692.\\[0pt]
[19] M. Matsui, Linear cryptanalysis method for DES cipher, in: Advances in Cryptology - EUROCRYPT'93, in: Lecture Notes in Comput. Sci., Springer-Verlag, 1994, pp. 386-397.\\[0pt]
[20] K. Nyberg, Differentially uniform mappings for cryptography, in: Advances in Cryptography, EUROCRYPT'93, in: Lecture Notes in Comput. Sci., vol. 765, Springer-Verlag, New York, 1994, pp. 55-64.\\[0pt]
[21] K. Nyberg, S-boxes and round functions with controllable linearity and differential uniformity, in: Proceedings of FSE 1994, in: Lecture Notes in Comput. Sci., vol. 1008, 1995, pp. 111-130.

\begin{itemize}
  \item 
\end{itemize}


\end{document}